\section{Library architecture: knowledge, computation, and verification}
\label{sec:framework}
\textsc{fin-skills} connects three reusable layers. Skills describe financial assumptions
and library conventions. Numerical interfaces execute methods against declared inputs.
Research guards return structured findings that a controller can require before accepting
an artifact. A successful algorithm call alone does not imply that the required audit ran.

\paragraph{Guarded research workflow.}
For a task and its available data, an agent selects a skill, inspects a numerical
interface, and proposes an artifact. A controller executes the relevant checks on that
artifact. Failures return actionable findings; missing required evidence remains
incomplete. The controller accepts only artifacts meeting its declared policy and saves
the associated execution records. This workflow is distinct from offering an agent
optional tool access and trusting its claim to have validated the output.

\begin{table*}[t]
\centering\small
\setlength{\tabcolsep}{5pt}
\begin{tabular}{p{0.16\textwidth}p{0.26\textwidth}p{0.25\textwidth}p{0.23\textwidth}}
\toprule
Component & Role in research workflow & Target question & Implementation and validation \\
\midrule
Financial skills & Route tasks and document dated conventions &
What assumptions and failure modes apply? & 129 skill documents \\
Numerical interfaces & Discover and execute quantitative methods &
Do selected computations agree with references? & 48 catalog entries; E3 tests eight primitives \\
Executable guards & Diagnose data, accounting, and validation defects &
Which required checks passed, failed, or lack inputs? & 36 registered guards; E1 exercises 13 \\
Execution records & Bind verification to the submitted artifacts &
Did the required checks actually run on this code and data? & Explicit pass, fail, and incomplete states \\
Optional extensions & Select context and retain verified experience &
How do auxiliary context and memory services affect tasks? & HiSTrim and persistent memory; separate studies \\
\bottomrule
\end{tabular}
\caption{Core financial library components and optional context and memory integrations.}
\label{tab:framework}
\end{table*}

\paragraph{Optional experience services.}
HiSTrim operates within a single session by evicting tokens from the KV cache during
generation. Persistent memory operates across sessions by storing parameter updates in an
external module. The two mechanisms operate on different timescales: context pruning
controls what the model attends to right now, whereas persistent memory retains summary
statistics of what happened in previous episodes. Neither service is required to run
the library's guards or reproduce its numerical primitives.

\paragraph{Grounded feedback updates.}
When an episode concludes, the agent records the market context, the chosen tool, and the
realized return. Only verified, settlement-cleared outcomes update memory. Negative returns
are retained as learning signals, while trades relying on look-ahead data or unadjusted
prices are rejected by guards to prevent memory corruption.
